\documentclass[11pt]{article}
\usepackage[a4paper,margin=30mm]{geometry}
\usepackage{amsmath,amssymb,amsthm,mathtools}
\usepackage{booktabs,longtable}
\usepackage[hidelinks]{hyperref}
\usepackage{microtype}

\newtheorem{theorem}{Theorem}
\newtheorem{lemma}{Lemma}
\newcommand{\C}{\mathbb C}
\newcommand{\R}{\mathbb R}
\newcommand{\PP}{\mathbb P}
\newcommand{\cE}{\mathcal E}
\newcommand{\tr}{\operatorname{tr}}

\title{An Exact Five-Dimensional Counterexample to Concavity of Determinantal-Process Entropy}
\author{Author placeholder}
\date{Internal draft --- 5 September 2026}

\begin{document}
\maketitle

\begin{abstract}
For a determinantal probability measure on a finite labelled set, Russell Lyons
conjectured that the Shannon entropy of the full subset-valued law is concave in
the positive-contraction kernel.  We give an exact counterexample on five
coordinates over a complex Hilbert space.  The midpoint is a rational real
symmetric kernel, the chord has a rational purely imaginary Hermitian direction,
and every claim is reduced to finite rational arithmetic.  Sylvester certificates
put the midpoint and both endpoints strictly between zero and the identity.  A
table gives all 32 atom probabilities at the midpoint and endpoints; every entry
is positive, every distribution sums to one, and the two endpoint distributions
agree entry by entry.  Directed rational enclosures for logarithms prove
\[
 H(K+hA)-H(K)>10^{-13}
 \quad\text{and}\quad
 \frac1{200}<D^2H(K)[A,A]<\frac1{190}.
\]
Thus the complex-Hilbert-space version of the conjecture is false.  The
real-symmetric restriction and the minimum possible counterexample dimension
remain open.
\end{abstract}

\paragraph{Draft status.}
This is a private mathematical note awaiting an independent paper-level review.
All load-bearing prose and proof organization introduced here are marked
\texttt{NEW\_REWRITE\_REQUIRES\_VERIFICATION} in \texttt{source\_map.md}.

\section{The conjecture and the exact conclusion}

Let $E=[n]$ be a labelled set and let $K$ be a Hermitian operator on
$\ell^2(E;\C)$ with $0\preceq K\preceq I$.  The determinantal probability
measure $\PP^K$ is characterized by
\begin{equation}\label{eq:inclusion}
 \PP^K(T\subseteq X)=\det K[T]\qquad(T\subseteq E),
\end{equation}
where $K[T]$ is the principal submatrix on the fixed coordinates $T$.  Write
$p_K(S)=\PP^K(X=S)$ for its $2^n$ atom probabilities and
\begin{equation}\label{eq:entropy}
 H(K)=-\sum_{S\subseteq E}p_K(S)\log p_K(S),
 \qquad 0\log0:=0.
\end{equation}
Conjecture~2.6 of Lyons~\cite{Lyons2014} asserts midpoint concavity of this
quantity on the convex set of positive contractions.  The source explicitly
allows real or complex Hilbert spaces, defines a positive contraction to be
self-adjoint with $0\preceq K\preceq I$, and uses the entropy of the complete
subset-valued distribution.  These points occur, respectively, on printed pages
138, 140, and 141; the coordinate DPP definition begins on printed page 137.

\begin{theorem}\label{thm:main}
On $E=[5]$, there exist strict positive contractions $K-hA$ and $K+hA$ on
$\ell^2(E;\C)$, with midpoint $K$, such that
\[
 H(K)<\frac{H(K-hA)+H(K+hA)}2.
\]
More precisely, the matrices in \eqref{eq:K}--\eqref{eq:B} with
$h=1/100000$ satisfy
\[
 H(K+hA)-H(K)>10^{-13},
 \qquad
 \frac1{200}<D^2H(K)[A,A]<\frac1{190}.
\]
Consequently the full complex-Hermitian version of Lyons's entropy-concavity
conjecture is false.
\end{theorem}

The assertion concerns the labelled-coordinate entropy in \eqref{eq:entropy}.
No invariance under an arbitrary unitary change of basis is used or asserted.

\section{The rational chord}

Define
\begin{equation}\label{eq:K}
K=\frac1{10000}
\begin{pmatrix}
2719&-3449&1009&-818&-2490\\
-3449&5840&-2889&-1513&1322\\
1009&-2889&2148&2505&1101\\
-818&-1513&2505&4698&3958\\
-2490&1322&1101&3958&4596
\end{pmatrix}
\end{equation}
and $A=iB$, where
\begin{equation}\label{eq:B}
B=\frac1{50}
\begin{pmatrix}
0&-3&-13&-16&4\\
3&0&-4&15&17\\
13&4&0&6&-15\\
16&-15&-6&0&1\\
-4&-17&15&-1&0
\end{pmatrix},
\qquad h=\frac1{100000}.
\end{equation}
Since $B^T=-B$, the matrix $A=iB$ is Hermitian.  The Pfaffian numerator of
the leading $4\times4$ block of $50B$ is
$(-3)(6)-(-13)(15)+(-16)(-4)=241$, so that block has determinant
$58081/6250000>0$.  An odd-dimensional real antisymmetric matrix has zero
determinant; hence $\operatorname{rank}A=\operatorname{rank}B=4$.

\section{Strict feasibility by Sylvester's criterion}

\begin{lemma}\label{lem:feasible}
The three kernels $K$, $K-hA$, and $K+hA$ satisfy $0<K_t<I$ strictly.
\end{lemma}

\begin{proof}
For each row below, the five entries are the leading principal minors in orders
$1,\ldots,5$:
\begin{align*}
K:\;&
\frac{2719}{10^4},\;
\frac{3983359}{10^8},\;
\frac{24692191}{10^{12}},\\[-2pt]
&\frac{3908717}{5\cdot10^{14}},\;
\frac{118205061}{10^{20}};\\
I-K:\;&
\frac{7281}{10^4},\;
\frac{18393359}{10^8},\;
\frac{59312197809}{10^{12}},\\[-2pt]
&\frac{35054548217}{5\cdot10^{14}},\;
\frac{1709337154939}{10^{20}};\\
K+hA:\;&
\frac{2719}{10^4},\;
\frac{995839749991}{25\cdot10^{12}},\;
\frac{123460921839}{5\cdot10^{15}},\\[-2pt]
&\frac{4883537625128308081}{625\cdot10^{24}},\;
\frac{7353066474454085031}{625\cdot10^{28}};\\
I-K-hA:\;&
\frac{7281}{10^4},\;
\frac{4598339749991}{25\cdot10^{12}},\;
\frac{296560989039361}{5\cdot10^{15}},\\[-2pt]
&\frac{43818177021217628308081}{625\cdot10^{24}},\;
\frac{106676909868459401834969}{625\cdot10^{28}}.
\end{align*}
Every displayed number is a positive rational number, and the imaginary part of
each determinant is exactly zero.  Sylvester's criterion proves $K>0$, $I-K>0$,
$K+hA>0$, and $I-K-hA>0$.  Complex conjugation gives
$K-hA=\overline{K+hA}$ and $I-K+hA=\overline{I-K-hA}$, so their corresponding
leading principal minors are the same.  Thus both remaining Hermitian matrices
are positive definite as well.
\end{proof}

\section{All atom probabilities}

For $S\subseteq[5]$, put $C=S^c$, let $I_C$ denote the diagonal coordinate
projection onto $C$, and set
\begin{equation}\label{eq:eventdet}
 M_S=K-I_C,
 \qquad
 p_S(t)=(-1)^{|C|}\det(M_S+t iB).
\end{equation}

\begin{lemma}\label{lem:eventformula}
For every real $t$ for which $K+tA$ is a positive contraction,
$p_S(t)=\PP^{K+tA}(X=S)$.
\end{lemma}

\begin{proof}
M\"obius inversion of \eqref{eq:inclusion} gives
\begin{equation}\label{eq:mobius}
 \PP^{K+tA}(X=S)=
 \sum_{T\supseteq S}(-1)^{|T\setminus S|}\det(K+tA)[T].
\end{equation}
Expanding $\det(K+tA-I_C)$ by the diagonal entries contributed by $-I_C$
leaves exactly the same principal minors, with the same signs after multiplication
by $(-1)^{|C|}$.  This proves \eqref{eq:eventdet}.
\end{proof}

Table~\ref{tab:events} lists the exact data for all 32 masks; a bit equal to one
means that the corresponding coordinate belongs to $S$.  All center numerators
$c_S$ and endpoint numerators $e_S$ are strictly positive, and direct integer
summation gives
\begin{equation}\label{eq:sums}
 \sum_Sc_S=10^{20},
 \qquad
 \sum_Se_S=625\cdot10^{28}.
\end{equation}
Thus the center and endpoint rows are positive probability distributions and each
sums exactly to one.  Moreover,
\[
 p_S(-h)=\overline{p_S(h)}=p_S(h),
\]
because the matrices in \eqref{eq:eventdet} are Hermitian and changing $h$ to
$-h$ complex-conjugates them.  Hence the two endpoint distributions agree for
each labelled subset, not merely after a permutation or in aggregate.

\section{The exact logarithm certificates}

We now describe the directed interval calculation used for both strict
inequalities.  It involves rational arithmetic only.  For a positive rational
$x$, write uniquely $x=2^k y$ with $1\le y<2$, and put
$z=(y-1)/(y+1)\in[0,1/3)$.  For $N=80$, define
\[
 L_N(y)=2\sum_{m=0}^{N-1}\frac{z^{2m+1}}{2m+1},
 \qquad
 R_N(y)=\frac{2z^{2N+1}}{(2N+1)(1-z^2)}.
\]
The positive-term expansion of $\operatorname{atanh}z$ gives the rational
enclosure
\begin{equation}\label{eq:loginterval}
 L_N(y)\le\log y\le L_N(y)+R_N(y).
\end{equation}
The same formula at $y=2$, $z=1/3$, encloses $\log2$.  Combining the intervals
for $\log y$ and $k\log2$ uses the lower endpoint of the $\log2$ interval when
$k\ge0$ and the upper endpoint when $k<0$; the opposite endpoints give the upper
bound.  Finally, in a sum $\sum_j a_j\log x_j$, the lower bound uses the lower
logarithm endpoint when $a_j>0$ and the upper endpoint when $a_j<0$.  The upper
bound uses the reverse choices.  This explicit sign rule is what makes the
certificates directed.

\subsection{The midpoint Hessian}

Since $M_S$ is real symmetric and $B$ is real antisymmetric,
$\tr(M_S^{-1}B)=0$.  The determinant derivative formula applied to
\eqref{eq:eventdet} therefore gives
\begin{equation}\label{eq:pderivatives}
 p'_S(0)=0,
 \qquad
 p''_S(0)=p_S(0)\tr(M_S^{-1}B M_S^{-1}B).
\end{equation}
The last column of Table~\ref{tab:events} lists the numerator $d_S$ of every
$p''_S(0)$ over the common denominator $312500000000000$.  Exact summation gives
$\sum_Sp''_S(0)=0$.  Twice differentiating \eqref{eq:entropy}, using
$p'_S(0)=0$, yields
\begin{equation}\label{eq:hessianlogs}
 D^2H(K)[A,A]
 =-\sum_Sp''_S(0)\log p_S(0)
 =\sum_Sq_S\log p_S(0),
 \qquad q_S=-p''_S(0).
\end{equation}
Applying the directed rule above to the 32 rational logarithms gives the strict
rational enclosure
\begin{equation}\label{eq:hessianinterval}
 \frac{5205948204140858303}{10^{21}}
 <D^2H(K)[A,A]<
 \frac{5205948204140858304}{10^{21}}.
\end{equation}
The lower endpoint exceeds $1/200$ by
$205948204140858303/10^{21}$, while $1/190$ exceeds the upper endpoint by
$16984126895682691/296875000000000000000$.  Both margins are positive, which
proves the Hessian claim without using its decimal approximation.

\subsection{The finite chord}

Use the 64 rational logarithmic terms in
\begin{equation}\label{eq:gaplogs}
 H(K+hA)-H(K)
 =-\sum_S p_S(h)\log p_S(h)
  +\sum_S p_S(0)\log p_S(0).
\end{equation}
The probabilities are exactly those in Table~\ref{tab:events}.  Applying
\eqref{eq:loginterval}, with each coefficient directed according to its sign,
gives
\begin{equation}\label{eq:gapinterval}
 \frac{233555253804221762}{10^{30}}
 <H(K+hA)-H(K)<
 \frac{233555253804221763}{10^{30}}.
\end{equation}
The lower endpoint exceeds $10^{-13}$ by
\[
 \frac{66777626902110881}
 {500000000000000000000000000000}>0.
\]
Because the two endpoint distributions agree term by term,
$H(K-hA)=H(K+hA)$.  Lemma~\ref{lem:feasible} and
\eqref{eq:gapinterval} now prove Theorem~\ref{thm:main}.

\section{Scope and compatibility}

The counterexample uses a real symmetric midpoint but a purely imaginary
rank-four Hermitian direction.  It therefore settles the complex-Hilbert-space
branch of the original conjecture, whose quantifiers include this chord.  It does
not settle a separately imposed restriction to real symmetric kernels.  It also
does not show that five is the least possible counterexample dimension.

This result is compatible with the independently verified fixed-complement-block
star-fibre theorem.  In that theorem the kernels have the form
\[
 \begin{pmatrix}a&v^*\\v&C\end{pmatrix}
\]
with the same Hermitian block $C$ along the fibre, so a difference within one
fibre is supported on one coordinate star and has rank at most two.  The present
direction has rank four and is not a difference inside such a fibre.  Coordinate
dependence also prevents an arbitrary unitary change of basis from turning this
observation into a claim about the same entropy functional.

The exact certificate proves a counterexample, but it does not establish absolute
priority or novelty.  The source audit found no earlier public proof or
counterexample in its bounded search as of 5 September 2026; that limited result
is not a systematic or exhaustive priority determination.

\section*{Data availability}

The directory \texttt{certificates/} contains the two frozen 32-row rational
tables and their JSON summaries.  Table~\ref{tab:events} is a typeset projection
of those fixed tables, not a floating-point recomputation.

\appendix
\section{The 32 exact atoms and second derivatives}
\begingroup\scriptsize
\begin{longtable}{@{}c r r r@{}}
\caption{Exact event data. For each mask $S$, the center probability is $c_S/10^{20}$, either endpoint probability is $e_S/(625\cdot10^{28})$, and $p_S''(0)=d_S/312500000000000$. Thus $q_S=-d_S/312500000000000$ in the Hessian formula.}\label{tab:events}\\
\toprule
mask & $c_S$ & $e_S$ & $d_S$\\
\midrule
\endfirsthead
\toprule
mask & $c_S$ & $e_S$ & $d_S$\\
\midrule
\endhead
\texttt{00000} & 1709337154939 & 106676909868459401834969 & -156662315231953\\
\texttt{00001} & 2656283230935061 & 166017744649677366641455031 & 42716236056953\\
\texttt{00010} & 6115970682105061 & 382248259084754513622575031 & 91453188204453\\
\texttt{00011} & 3971055587349804939 & 248190974216992099660334134969 & 7629290970547\\
\texttt{00100} & 3263125734325061 & 203945390776741244180575031 & 32381424934453\\
\texttt{00101} & 4816014638897584939 & 301000914895668312929776134969 & -35430745759453\\
\texttt{00110} & 4184669109046414939 & 261541819326161410782795014969 & 10760477093047\\
\texttt{00111} & 2467575721675061 & 154223489507135043248275031 & 6902443731953\\
\texttt{01000} & 8434437215975061 & 527152398460111688723165031 & 72461670379453\\
\texttt{01001} & 12094797748615934939 & 755924859218562692485233544969 & -69933241204453\\
\texttt{01010} & 25132478169064764939 & 1570779885473979040338252424969 & -92568768351953\\
\texttt{01011} & 4467094503325061 & 279193455111755487790865031 & 48653939176953\\
\texttt{01100} & 3805671766912544939 & 237854485447237453607694424969 & 15203394918047\\
\texttt{01101} & 2116290055545061 & 132268155485142218348865031 & 27013575906953\\
\texttt{01110} & 5789712006715061 & 361857023390619365329985031 & 22970928054453\\
\texttt{01111} & 781625194939 & 48828023184808626724969 & -23551498879453\\
\texttt{10000} & 7009200306245061 & 438075093302307823678975031 & 74161991514453\\
\texttt{10001} & 6285970421125664939 & 392873151332436946350277734969 & 12082887660547\\
\texttt{10010} & 25078103510574494939 & 1567381469304727819203296614969 & -106178114486953\\
\texttt{10011} & 5366317393595061 & 335394862896526622746675031 & 25796835311953\\
\texttt{10100} & 8650505306722274939 & 540656581615572207472738614969 & -54569976216953\\
\texttt{10101} & 1730314645815061 & 108144694433938353304675031 & 29070497041953\\
\texttt{10110} & 5070067596985061 & 316879269464890500285795031 & 44653324189453\\
\texttt{10111} & 861634924939 & 53827415363673670914969 & -24767445014453\\
\texttt{11000} & 5915774434040624939 & 369735902130528962028196024969 & 2989903338047\\
\texttt{11001} & 3354766127465061 & 209672901149433797847265031 & 18182867486953\\
\texttt{11010} & 4413568778635061 & 275848120757135944828385031 & 72092444634453\\
\texttt{11011} & 781653274939 & 48810450868229128324969 & -42878815459453\\
\texttt{11100} & 2699019730855061 & 168688744962347675386385031 & 11783906364453\\
\texttt{11101} & 537301054939 & 33559863856498570324969 & -21452077189453\\
\texttt{11110} & 892249884939 & 55724684329351589204969 & -40933479336953\\
\texttt{11111} & 118205061 & 7353066474454085031 & -34749838047\\
\bottomrule
\end{longtable}

\endgroup

\bibliographystyle{plain}
\bibliography{references}

\end{document}
